% Part: applied-modal-logics
% Chapter: temporal-logic
% Section: language-epistemic-logic

\documentclass[../../../include/open-logic-section]{subfiles}

\begin{document}

\olfileid[hi]{aml}{tl}{sem}

\olsection{कालिक तर्कशास्त्र का अर्थविज्ञान}

\begin{defn}
कालिक तर्कशास्त्र की मूल भाषा में निम्नलिखित होते हैं:
\begin{enumerate}
  \tagitem{prvFalse}{!!{falsity} का प्रतिज्ञप्तिक अचर~$\lfalse$।}{}
  \tagitem{prvTrue}{!!{truth} का प्रतिज्ञप्तिक अचर~$\ltrue$।}{}
  \item !!{propositional variable}ों का !!{denumerable} समुच्चय: $\Obj
    p_0$, $\Obj p_1$, $\Obj p_2$, \dots
  \item प्रतिज्ञप्तिक संयोजक: \startycommalist
  \iftag{prvNot}{\ycomma $\lnot$ (निषेध)}{}%
  \iftag{prvAnd}{\ycomma $\land$ (संयोजन)}{}%
  \iftag{prvOr}{\ycomma $\lor$ (वियोजन)}{}%
  \iftag{prvIf}{\ycomma $\lif$ (!!{conditional})}{}%
  \iftag{prvIff}{\ycomma $\liff$ (!!{biconditional})}{}।
  \item अतीत के संकारक $\Ptemp$ और $\Htemp$।
  \item भविष्य के संकारक $\Ftemp$ और $\Gtemp$।
\end{enumerate}
\end{defn}

आगे हम अन्य प्रकार के मोडल संकारक जोड़ने की संभावना पर चर्चा करेंगे।

\begin{defn}
कालिक भाषा की \emph{!!{formula}ें} आगमनात्मक रूप से इस प्रकार परिभाषित होती
हैं:
\begin{enumerate}
\tagitem{prvFalse}{$\lfalse$ एक परमाण्विक !!{formula} है।}{}

\tagitem{prvTrue}{$\ltrue$ एक परमाण्विक !!{formula} है।}{}

\item प्रत्येक प्रतिज्ञप्तिक चर $\Obj p_i$ एक (परमाण्विक) !!{formula} है।

\tagitem{prvNot}{यदि $!A$ एक !!{formula} है, तो $\lnot !A$ एक
  !!{formula} है।}{}

\tagitem{prvAnd}{यदि $!A$ और $!B$ !!{formula} हैं, तो $(!A \land
  !B)$ एक !!{formula} है।}{}

\tagitem{prvOr}{यदि $!A$ और $!B$ !!{formula} हैं, तो $(!A \lor !B)$
  एक !!{formula} है।}{}

\tagitem{prvIf}{यदि $!A$ और $!B$ !!{formula} हैं, तो $(!A \lif !B)$
  एक !!{formula} है।}{}

\tagitem{prvIff}{यदि $!A$ और $!B$ !!{formula} हैं, तो $(!A \liff !B)$
  एक !!{formula} है।}{}

\item यदि $!A$ एक !!{formula} है, तो $\Ptemp !A$, $\Htemp !A$, $F !A$ और
  $\Gtemp !A$ सभी !!{formula} हैं।

\tagitem{limitClause}{इसके अतिरिक्त कुछ भी !!{formula} नहीं है।}{}
\end{enumerate}
\end{defn}

अन्य प्रकार के आशयात्मक तर्कशास्त्रों की तरह कालिक तर्कशास्त्रों का अर्थविज्ञान
भी संबंधात्मक प्रतिरूपों के द्वारा दिया जाता है।

\begin{defn}
  कालिक भाषा का \emph{प्रतिरूप} एक त्रिक
  $\mModel{M} = \tuple{T, \prec, V}$ है, जहाँ
  \begin{enumerate}
  \item $T$ एक अरिक्त समुच्चय है, जिसके अवयव समय-बिंदु माने जाते हैं।
  \item $\prec$, $T$ पर एक द्विपदी संबंध है।
  \item $V$ ऐसा फलन है जो प्रत्येक !!{propositional variable}~$p$ को समय-बिंदुओं
    का समुच्चय $V(p)$ देता है।
  \end{enumerate}
  जब $t \prec t'$ हो, तो हम कहते हैं कि $t$, $t'$ से \emph{पहले} है।
  जब $t \in V(p)$ हो, तो हम कहते हैं कि $p$, $t$ पर \emph{सत्य} है।
\end{defn}

फिलहाल हम अपने पूर्वता-संबंध $\prec$ पर कोई शर्त नहीं लगाते। इसका अर्थ है
कि अभी हमारे कालिक प्रतिरूपों की संरचना पर कोई प्रतिबंध नहीं है; अतः हमारे
पास ऐसे प्रतिरूप हो सकते हैं जिनमें समय रैखिक, शाखित, चक्रीय अथवा किसी भी
अन्य संरचना वाला हो।

सामान्य मोडल तर्कशास्त्र की ही तरह प्रत्येक कालिक प्रतिरूप निर्धारित करता
है कि उसमें कौन-से !!{formula} किन बिंदुओं पर सत्य माने जाते हैं। संबंधात्मक
अर्थविज्ञान की मूल धारणा के लिए हम वही संकेत-पद्धति---``प्रतिरूप
$\mModel{M}$, !!{formula}~$!A$ को बिंदु~$t$ पर सत्य बनाता है''---अपनाते हैं।
यह संबंध आगमनात्मक रूप से परिभाषित होता है और सभी गैर-मोडल संकारकों के लिए
सामान्य मोडल मामले के समान है।

\begin{defn}\ollabel{defn:tmodels}
  किसी $\mModel M$ में \emph{!!{formula}~$!A$ की $t$ पर सत्यता}, जिसे
  $\mSat{M}{!A}[t]$ से लिखते हैं, आगमनात्मक रूप से इस प्रकार परिभाषित है:
  \begin{enumerate}
  \tagitem{prvFalse}{\indcase{!A}{\lfalse}{$\mSat{M}{\lfalse}[t]$ कभी नहीं।}}{}
  \tagitem{prvTrue}{\indcase{!A}{\ltrue}{$\mSat{M}{\ltrue}[t]$ सदैव।}}{}
  \item $\mSat{M}{p}[t]$ तभी जब $t \in V(p)$।
  \tagitem{prvNot}{\indcase{!A}{\lnot !B}{$\mSat{M}{\indfrm}[t]$ तभी जब
    $\mSat/{M}{!B}[t]$।}}{}
  \tagitem{prvAnd}{\indcase{!A}{(!B \land !C)}{$\mSat{M}{\indfrm}[t]$ तभी जब
    $\mSat{M}{!B}[t]$ और $\mSat{M}{!C}[t]$।}}{}
  \tagitem{prvOr}{\indcase{!A}{(!B \lor !C)}{$\mSat{M}{\indfrm}[t]$ तभी जब
    $\mSat{M}{!B}[t]$ या $\mSat{M}{!C}[t]$} (या दोनों)।}{}
  \tagitem{prvIf}{\indcase{!A}{(!B \lif !C)}{$\mSat{M}{\indfrm}[t]$ तभी जब
    $\mSat/{M}{!B}[t]$ या $\mSat{M}{!C}[t]$।}}{}
  \tagitem{prvIff}{\indcase{!A}{(!B \liff !C)}{$\mSat{M}{\indfrm}[t]$ तभी जब
    या तो $\mSat{M}{!B}[t]$ और $\mSat{M}{!C}[t]$ दोनों, या
    न $\mSat{M}{!B}[t]$ न $\mSat{M}{!C}[t]$।}}{}
  \item\ollabel{defn:sub:mmodels-p}
    \indcase{!A}{\Ptemp  !B}{$\mSat{M}{\indfrm}[t]$ तभी जब
    $t' \prec t$ वाले किसी $t' \in T$ के लिए $\mSat{M}{!B}[t']$।}
\item\ollabel{defn:sub:mmodels-h}
    \indcase{!A}{\Htemp  !B}{$\mSat{M}{\indfrm}[t]$ तभी जब
    $t' \prec t$ वाले प्रत्येक $t' \in T$ के लिए $\mSat{M}{!B}[t']$।}
   \item\ollabel{defn:sub:mmodels-f}
    \indcase{!A}{\Ftemp  !B}{$\mSat{M}{\indfrm}[t]$ तभी जब
    $t \prec t'$ वाले किसी $t' \in T$ के लिए $\mSat{M}{!B}[t']$।}
\item\ollabel{defn:sub:mmodels-g}
    \indcase{!A}{\Gtemp  !B}{$\mSat{M}{\indfrm}[t]$ तभी जब
    $t \prec t'$ वाले प्रत्येक $t' \in T$ के लिए $\mSat{M}{!B}[t']$।}
  \end{enumerate} 
\end{defn}

अर्थविज्ञान से देखा जा सकता है कि संकारक $\Ptemp$ और~$\Htemp$ एक-दूसरे के
द्वैत हैं, और इसी प्रकार $\Ftemp$ तथा $\Gtemp$ भी। अतः हम $\Htemp !A$ को
$\lnot \Ptemp \lnot !A$ से परिभाषित कर सकते थे, और $\Gtemp$ को भी
$\Ftemp$ के द्वारा इसी तरह परिभाषित कर सकते थे।

\end{document}
